% Machine-bound publication statement for A3_KCLASS_ENDPOINT_DP_THEOREM_V1.
% Normative authority: PUBLICATION_MANIFEST.json and the admitted receipt pinned there.
\section{Main result}

Let $\mathbb F$ be a finite field and let $U_1,\ldots,U_k$ be the distinct geometric subspaces occurring in a finite subspace arrangement.  Class $j$ occurs with positive multiplicity $a_j$.  Write $s=|\{j:a_j=1\}|$ and $r=|\{j:a_j\ge 2\}|$, so $s+r=k$.

\begin{theorem}[Endpoint-compressed ternary dynamic program]
For the admitted repeated-geometric-class subspace-arrangement model, replacing each singleton class by one atomic event and each repeated class by its first and last occurrence preserves the optimum path-width.  The resulting endpoint-state directed acyclic graph has exactly
\[
  2^s3^r
\]
admissible states and exactly
\[
  s2^{s-1}3^r+2r2^s3^{r-1}
\]
directed transitions, with a term interpreted as zero when its corresponding class type is absent.

Let
\[
  \rho(X)=\dim\!\left(\sum_{j\in X}U_j\right),\qquad X\subseteq\{1,\ldots,k\}.
\]
For a state whose finished classes are $P$ and whose started classes are $S$, the cut width is
\[
  \lambda(P,S)=\rho(S)+\rho(K\setminus P)-\rho(K).
\]
Consequently, after preprocessing the $2^k$ subset ranks, each state-width query requires constant lookup/arithmetic work and the bottleneck dynamic program computes the exact optimum with combinatorial work
\[
  O(k3^k).
\]
\end{theorem}

\paragraph{Claim boundary.}
This theorem concerns the stated finite-field repeated-geometric-class subspace-arrangement domain.  It does not alter the known complexity status of general matroid path-width, does not establish a polynomial-time algorithm when $k$ is unbounded, and makes no claim concerning $P$ versus $NP$.
